
Les travaux de \cite{romain25} s'inscrivent dans une approche de calcul à zones couplant conduction et rayonnement, où chaque surface est caractérisée par une unique température moyenne. Cette approche, bien qu'efficace, ne permet pas de rendre compte des écarts de température au sein d'une même surface, écarts pourtant susceptibles d'être significatifs dans des configurations à forte hétérogénéité de puissance.

Cette étude se place dans une perspective différente et complémentaire : il ne s'agit plus de raffiner un modèle à zones, mais d'exploiter directement les capacités de la simulation CFD 3D pour accéder, en plus de la température moyenne $T_{moy}$, aux températures extrémales $T_{max}$ et $T_{min}$ d'une géométrie représentative d'un assemblage de crayons combustible. L'objectif de ce chapitre est ainsi de démontrer la faisabilité d'une étude 3D exigeante sur une telle configuration, de quantifier le biais introduit par la discrétisation géométrique des surfaces sur ces trois grandeurs et d'estimer l'écart entre le point chaud et la température moyenne.